% ============================================================
% ANHANG BCLM-LL-MPS: MULTISPEZIES-VALIDIERUNGSPROTOKOLL
% FÜR DAS LANGLEBIGKEITS-KOORDINATIONSDESIGN
% ============================================================

\documentclass[12pt,a4paper]{article}

% Für XeLaTeX und deutsche Sprache (wie in vorherigen Übersetzungen)
\usepackage{fontspec}
\usepackage{polyglossia}
\setmainlanguage{german}

\usepackage{amsmath,amssymb,amsthm,mathtools}
\usepackage{geometry}
\geometry{left=1.5cm,right=1.5cm,top=1.5cm,bottom=2.0cm}
\usepackage{booktabs}
\usepackage{longtable}
\usepackage{array}
\usepackage{hyperref}
\usepackage{url}
\usepackage{xcolor}
\usepackage{titlesec}
\usepackage{fancyhdr}
\usepackage{enumitem}
\usepackage{makecell}
\usepackage{float}

\titleformat{\section}{\Large\bfseries}{\thesection}{1em}{}
\titleformat{\subsection}{\large\bfseries}{\thesubsection}{1em}{}

\hypersetup{colorlinks=true,linkcolor=blue,citecolor=blue,urlcolor=blue}

% --- YUCT fundamentale Konstanten ---
\newcommand{\REL}{\mathcal{K}}          % relative Koordinationseffizienz
\newcommand{\ABS}{K_{\mathrm{eff}}}     % absolute
\newcommand{\kappac}{\kappa_c}
\newcommand{\betaf}{\beta}
\newcommand{\Sodd}{S_{\mathrm{odd}}}
\newcommand{\Seven}{S_{\mathrm{even}}}

\begin{document}
	
	\begin{titlepage}
		\centering
		\vspace*{2cm}
		{\Huge\bfseries Anhang BCLM-LL-MPS\\[0.3cm]
			Multispezies-Validierungsprotokoll\\für das Langlebigkeits-Koordinationsdesign\par}
		\vspace{0.5cm}
		{\Large\bfseries Ein paralleles Experiment an fünf Säugetierarten\\zur Bestätigung der YUCT-Koordinationstheorie des Alterns\par}
		\vspace{1.5cm}
		{\large A.\,V.\,Yakushev\par}
		{\normalsize YUCT-Forschung, \url{https://yuct.org/}\par}
		\vspace{1.0cm}
		{\normalsize Juni 2026 (v1.2 – Multispezies-Konsistenz)}\par
		\vfill
		\begin{abstract}
			\noindent
			Der YUCT-biologische Koordinationslagrange (BCLM) ermöglicht die Vorhersage der Koordinationseffizienz \(K_{\mathrm{eff}}\) für jedes Protein oder jeden Signalweg allein aus Sequenzdaten. Um die Universalität der Koordinationstheorie des Alterns zu demonstrieren, stellen wir ein paralleles experimentelles Design für fünf weit verbreitete LaborSäugetiere vor: Maus (\textit{Mus musculus}), Ratte (\textit{Rattus norvegicus}), Schwein (\textit{Sus scrofa}), Javaneraffe (\textit{Macaca fascicularis}) und Hund (\textit{Canis familiaris}). Für jede Art identifizieren wir die orthologen Zielgene, berechnen ihre natives und optimierten relativen Koordinationseffizienzen \(\REL\) mit der BCLM-Pipeline und leiten quantitative Vorhersagen für Mutationsrate, oxidativen Stress, Proteinaggregation, Migrationsgeschwindigkeit und replikative Lebensspanne direkt aus dem universellen Fehlergesetz \(\varepsilon = \kappa_c \alpha \REL^{-\beta}\) (\(\beta=2/3\), \(\kappa_c=1/3\)) ab. Eine Vergleichstabelle enthält die entsprechenden Werte für menschliche Zellen (aus dem begleitenden Protokoll BCLM‑LL). Das Experiment stützt sich ausschließlich auf kommerziell erhältliche Zelllinien, öffentliche Sequenzdatenbanken und Standard-Zellbiologie-Assays. Ein kontrolliertes Regressionsmodul (CCR) bestätigt die Reversibilität des Alterungsphänotyps über die Spezies hinweg. Alle Vorhersagen sind vorregistriert und falsifizierbar. Das Protokoll verzichtet bewusst auf \textit{in vivo}-Verabreichungsmethoden; ethische Aspekte werden in Anhang~\(\Sigma\) diskutiert.
		\end{abstract}
	\end{titlepage}
	
	\tableofcontents
	\newpage
	
	\section{Einführung}
	\label{sec:intro}
	Das universelle Fehlerskalierungsgesetz von YUCT,
	\begin{equation}
		\varepsilon = \kappa_c \, \alpha \, \REL^{-\beta}, \qquad \beta = \frac{2}{3},\quad \kappa_c = \frac{1}{3},
		\label{eq:error-law}
	\end{equation}
	gilt für jedes koordinierte System, unabhängig von Skala oder Substrat. In einem zellulären Kontext quantifiziert die relative Effizienz \(\REL\) die Genauigkeit molekularer Maschinen; ihr altersabhängiger Rückgang beschleunigt die Schadensakkumulation und löst schließlich Seneszenz aus.
	
	Der biologische Koordinationslagrange (BCLM, Anhang BCLM) liefert eine einheitliche Methode zur Berechnung von \(\REL\) aus Sequenzdaten und zum Entwurf codonoptimierter Konstrukte, die \(\REL\) auf juveniles Niveau zurückführen. Das begleitende Protokoll BCLM‑LL beschrieb ein vollständiges Experiment für humane Kardiomyozyten. Hier erweitern wir das Design auf fünf weitere Säugetierarten – Maus, Ratte, Schwein, Javaneraffe und Hund – jede mit vollständig annotierten Genomen, charakterisierter Codonnutzung und kommerziell erhältlichen Primärzelllinien. Ziel ist es, die Universalität des YUCT-Alterungsmechanismus zu testen: Wenn dieselbe quantitative Beziehung zwischen \(\REL\) und zellulärer Gesundheitsspanne für alle Arten gilt, wird die Koordinationstheorie des Alterns stark gestützt.
	
	\subsection{Vorhandene experimentelle Unterstützung}
	Die molekularen Ziele des Langlebigkeitsdesigns werden einzeln durch veröffentlichte Arbeiten gestützt:
	\begin{itemize}
		\item Eine hypergenaue Mutation im Ribosom verlängert die Lebensspanne in Hefe, Würmern und Fliegen \cite{Bjedov2021}.
		\item Langlebige Säugetiere zeigen eine Verarmung an hypermutablen Codons in Proteostasegenen \cite{Kerekes2025}.
		\item Überexpression mitochondrialer Chaperone reduziert oxidativen Schaden und verlängert die Lebensspanne von \textit{Drosophila} \cite{Morrow2004}.
		\item Verbesserte Proteostase verzögert aggregationsvermittelte Proteotoxizität \cite{Walker2003}.
		\item Altersbedingte Dysregulation der Integrin-Signalübertragung trägt zur zellulären Seneszenz bei \cite{Takeda2019}.
	\end{itemize}
	Keine dieser Studien führte jedoch einen einzigen quantitativen Koordinationsparameter ein. BCLM tut genau dies, und das untenstehende Multispezies-Protokoll liefert einen entscheidenden Test.
	
	\section{Zielarten und Zelllinien}
	\label{sec:species}
	Fünf Arten wurden ausgewählt, um einen weiten phylogenetischen Bereich abzudecken und gleichzeitig praktische Vorteile für die Zellkultur zu erhalten. Für jede sind primäre Fibroblasten der am leichtesten zugängliche normale Zelltyp, der über viele Passagen ohne Transformation vermehrbar ist.
	
	\begin{table}[H]
		\centering
		\caption{Die fünf Arten und ihre primären Zelllinien.}
		\label{tab:species}
		\begin{tabular}{@{}lllll@{}}
			\toprule
			\textbf{Art} & \textbf{Trivialname} & \textbf{Zelltyp} & \textbf{Quelle} \\
			\midrule
			\textit{Mus musculus} (C57BL/6) & Hausmaus & Embryonale Fibroblasten (MEF) & ATCC CRL‑2214 \\
			\textit{Rattus norvegicus} (F344) & Wanderratte & Dermale Fibroblasten & Cell Biolabs RAT‑DF \\
			\textit{Sus scrofa} (Landrasse) & Hausschwein & Dermale Fibroblasten & Cell Applications 12405 \\
			\textit{Macaca fascicularis} & Javaneraffe & Dermale Fibroblasten & Coriell AG08333 \\
			\textit{Canis familiaris} (Beagle) & Haushund & Dermale Fibroblasten & Coriell AG07090 \\
			\bottomrule
		\end{tabular}
	\end{table}
	
	\section{Orthologe Zielgene und ihre \(\REL\)}
	\label{sec:genes}
	Für jede Art identifizierten wir die Orthologe der menschlichen Proteine, die im Langlebigkeitsdesign (BCLM‑LL) anvisiert werden. Die Aminosäuresequenzen sind hoch konserviert; folglich ist der Aminosäurebeitrag zu \(K_{\mathrm{eff}}\) über die Arten hinweg im Wesentlichen identisch. Der Hauptunterschied kommt von der Codonoptimierung, da der Satz optimaler Codons von der artspezifischen Codonnutzungstabelle abhängt (Kazusa-Datenbank \cite{codon_usage_db}). Für jedes Gen berechneten wir zwei relative Koordinationseffizienzen: \(\REL_{\mathrm{WT}}\) (native Kodierungssequenz) und \(\REL_{\mathrm{LL}}\) (codonoptimierte Sequenz). Die Berechnung folgte den in BCLM, Abschnitt~2 beschriebenen Regeln.
	
	Tabelle~\ref{tab:targets} listet die Orthologe und die vorhergesagten \(\REL\)-Werte für die menschliche Referenz (aus BCLM‑LL) und für die fünf Tierarten auf. Alle Werte wurden mit dem Skript `bclm\_multi\_species.py` erzeugt, das im YPSDC-Repository verfügbar ist.
	
	\begin{table}[H]
		\centering
		\caption{Orthologe Zielgene und ihre vorhergesagten \(\REL\).}
		\label{tab:targets}
		{\footnotesize
			\begin{tabular}{@{}llcccccc@{}}
				\toprule
				\textbf{Ebene} & \textbf{Menschliches Gen} & \textbf{Maus} & \textbf{Ratte} & \textbf{Schwein} & \textbf{Javaneraffe} & \textbf{Hund} & \textbf{Fehlerfaktor\footnotemark[1]} \\
				\midrule
				1 & BRCA1 & 0,62→0,78 & 0,63→0,79 & 0,62→0,78 & 0,62→0,78 & 0,85 \\
				& PARP1 & 0,60→0,77 & 0,61→0,78 & 0,60→0,77 & 0,60→0,77 & 0,86 \\
				& MLH1  & 0,59→0,76 & 0,60→0,77 & 0,59→0,76 & 0,59→0,76 & 0,87 \\
				2 & NDUFS1 & 0,55→0,72 & 0,56→0,73 & 0,55→0,72 & 0,55→0,72 & 0,80 \\
				& NDUFV1 & 0,56→0,73 & 0,57→0,74 & 0,56→0,73 & 0,56→0,73 & 0,80 \\
				& NDUFA9 & 0,54→0,70 & 0,55→0,71 & 0,54→0,70 & 0,54→0,70 & 0,81 \\
				3 & HSPA1A & 0,52→0,69 & 0,53→0,70 & 0,52→0,69 & 0,52→0,69 & 0,78 \\
				& DNAJB1 & 0,50→0,68 & 0,51→0,69 & 0,50→0,68 & 0,50→0,68 & 0,79 \\
				& HSPA9  & 0,53→0,70 & 0,54→0,71 & 0,53→0,70 & 0,53→0,70 & 0,78 \\
				4 & ITGAV  & 0,58→0,74 & 0,59→0,75 & 0,58→0,74 & 0,58→0,74 & 0,83 \\
				& ITGB3  & 0,57→0,73 & 0,58→0,74 & 0,57→0,73 & 0,57→0,73 & 0,83 \\
				& FN1    & 0,60→0,76 & 0,61→0,77 & 0,60→0,76 & 0,60→0,76 & 0,84 \\
				\bottomrule
		\end{tabular}}
		\footnotetext[1]{Fehlerfaktor \(\varepsilon_{\mathrm{LL}}/\varepsilon_{\mathrm{WT}} = (\REL_{\mathrm{LL}}/\REL_{\mathrm{WT}})^{-2/3}\). Aufgrund der sehr hohen Sequenzkonservierung variiert der Faktor zwischen den Arten um \(<0,5\%\); die aufgeführten Werte sind die Mittelwerte.}
	\end{table}
	
	\section{Vorhergesagte phänotypische Veränderungen}
	\label{sec:predictions}
	Für jede Ebene wird die Gesamtverbesserung als geometrisches Mittel der einzelnen Proteinerrorfaktoren erhalten, da die anvisierten Komponenten in überlappenden Signalwegen wirken. Die resultierenden ebenenspezifischen Fehlerverhältnisse sind:
	
	\begin{itemize}
		\item \textbf{Ebene 1 – Genomstabilität:} \(\mu_{\mathrm{LL}}/\mu_{\mathrm{WT}} = \sqrt[3]{0,85\times0,86\times0,87} \approx 0,86\).
		\item \textbf{Ebene 2 – Mitochondriale ROS:} \(\mathrm{ROS}_{\mathrm{LL}}/\mathrm{ROS}_{\mathrm{WT}} = \sqrt[3]{0,80\times0,80\times0,81} \approx 0,80\).
		\item \textbf{Ebene 3 – Proteinaggregation:} \(\mathrm{Agg}_{\mathrm{LL}}/\mathrm{Agg}_{\mathrm{WT}} = \sqrt[3]{0,78\times0,79\times0,78} \approx 0,78\).
		\item \textbf{Ebene 4 – ECM-Adhäsionsfehler:} \(\mathrm{Err}_{\mathrm{LL}}/\mathrm{Err}_{\mathrm{WT}} = \sqrt[3]{0,83\times0,83\times0,84} \approx 0,83\). Da die Migrationsgeschwindigkeit umgekehrt zu Adhäsionsfehlern skaliert, gilt \(v_{\mathrm{LL}}/v_{\mathrm{WT}} \approx 1/0,83 = 1,20\).
	\end{itemize}
	
	Unter der Annahme, dass die vier Ebenen unabhängig voneinander versagen, ist die gesamte tödliche Fehlerwahrscheinlichkeit das Produkt der vier Verhältnisse:
	\[
	\frac{\varepsilon_{\mathrm{tödlich,LL}}}{\varepsilon_{\mathrm{tödlich,WT}}} \approx 0,86 \times 0,80 \times 0,78 \times 0,83 \approx 0,45.
	\]
	Unter der Hypothese, dass die replikative Lebensspanne (gemessen durch kumulative Populationsverdopplungen) umgekehrt proportional zur tödlichen Fehlerwahrscheinlichkeit ist, beträgt die erwartete Lebensspannenverlängerung
	\[
	\frac{\mathrm{PD}_{\mathrm{LL}}}{\mathrm{PD}_{\mathrm{WT}}} \approx \frac{1}{0,45} \approx 2,2.
	\]
	Dies ist eine obere Grenze; Übersprechen zwischen den Ebenen und verbleibende Schadenswege werden den tatsächlichen Gewinn reduzieren. Eine konservative Schätzung, die für die vorregistrierte Vorhersage übernommen wird, setzt den Lebensspannenverlängerungsfaktor für alle Arten in den Bereich \textbf{1,5–2,0}.
	
	Da die Aminosäuresequenzen nahezu identisch sind, sind diese Vorhersagen für alle fünf Tierarten gleich und stimmen mit der menschlichen Vorhersage aus BCLM‑LL überein.
	
	\section{Experimentelles Protokoll}
	\label{sec:protocol}
	
	\subsection{Genetische Konstrukte}
	Für jede Art werden zwei synthetische Varianten jedes Zielgens hergestellt:
	\begin{itemize}
		\item \textbf{WT‑OE}: die native Kodierungssequenz, kloniert in einen Doxycyclin-induzierbaren lentiviralen Vektor (pLV‑TRE‑Tight‑IRES‑EGFP).
		\item \textbf{LL‑OPT}: die BCLM-optimierte Sequenz (codonoptimiert für die spezifische Wirtsart) im selben Vektor.
	\end{itemize}
	Eine Scrambled-Kontrolle (SCR) wird für BRCA1 durch zufällige Permutation synonymer Codons erzeugt.
	
	\subsection{Kontrollen}
	Für jede Art werden drei Kontrollen parallel durchgeführt:
	\begin{enumerate}
		\item Leervektor (EV).
		\item WT‑OE (angepasste Proteinniveaus).
		\item SCR (Spezifitätskontrolle).
	\end{enumerate}
	
	\subsection{Messplan}
	Dieselben fünf Assays wie in BCLM‑LL werden für jede Art durchgeführt.
	
	\begin{table}[H]
		\centering
		\caption{Experimenteller Zeitplan.}
		\label{tab:schedule}
		\begin{tabular}{@{}llll@{}}
			\toprule
			\textbf{Tag} & \textbf{Assay} & \textbf{Ebene} & \textbf{Ergebnis} \\
			\midrule
			0–7   & Transduktion, Selektion & --- & EGFP \\
			7–14  & Doxycyclin-Induktion & Alle & Proteinniveaus \\
			14–21 & HPRT-Mutationsassay & 1 & 6‑TG\(^{\mathrm{R}}\)-Kolonien \\
			14–21 & MitoSOX, TMRE & 2 & Durchflusszytometrie \\
			21–28 & HTT‑Q72-Einschlussassay & 3 & Fluoreszenzmikroskopie \\
			21–28 & Kratzwundassay & 4 & Wundschlussgeschwindigkeit \\
			28–56 & Serielle Passagierung & Alle & Seneszenz-\(\beta\)-Gal, Telomer-qPCR \\
			\bottomrule
		\end{tabular}
	\end{table}
	
	\subsection{Bestimmung der replikativen Lebensspanne}
	Die Zellen werden bei 80\% Konfluenz passagiert, und die kumulativen Populationsverdopplungen (CPD) werden aufgezeichnet. Der primäre Endpunkt ist die CPD bei Wachstumsarrest. Für BCLM-optimierte Zellen wird eine \(1,5\)–\(2,0\)-fache CPD im Vergleich zu WT‑OE-Kontrollen vorhergesagt.
	
	\section{Speziesübergreifende kontrollierte Koordinations-Regression (CCR)}
	\label{sec:CCR}
	Die Reversibilität des Alterungsphänotyps wird in jeder Art getestet, indem ein optimiertes Gen pro Ebene transient unterdrückt wird.
	
	\begin{table}[H]
		\centering
		\caption{CCR-Interventionen für die fünf Arten.}
		\label{tab:CCR}
		\begin{tabular}{@{}llll@{}}
			\toprule
			\textbf{Ebene} & \textbf{Zielgen} & \textbf{Methode} & \textbf{Erwarteter Effekt} \\
			\midrule
			1 & BRCA1\_LL & Dox-induzierbare shRNA & \(\REL \to 0,62\); Mutationsrate steigt um Faktor \(1,86\) \\
			2 & NDUFS1\_LL & CRISPR-Baseneditierung (I182F) & \(\REL \to 0,55\); ROS steigt um Faktor \(1,80\) \\
			3 & HSPA1A\_LL & Tet‑CRISPRi & \(\REL \to 0,52\); Aggregate steigen um Faktor \(1,78\) \\
			4 & ITGB3\_LL & siRNA & \(\REL \to 0,57\); Migrationsgeschwindigkeit fällt um Faktor \(1,20\) \\
			\bottomrule
		\end{tabular}
	\end{table}
	
	Nach 72 h Unterdrückung müssen sich die Biomarker in Richtung des Wildtyp-Phänotyps verschieben. Nach Auswaschung oder siRNA-Verdünnung sollten die Zellen innerhalb von 48–72 h in den Langlebigkeitszustand zurückkehren. Die quantitative Vorhersage für jeden Biomarker wird direkt aus den \(\REL\)-Werten in Tabelle~\ref{tab:targets} und dem universellen Fehlergesetz abgeleitet.
	
	\section{Falsifizierbarkeit}
	\label{sec:falsifiability}
	Alle Vorhersagen in Tabelle~\ref{tab:predictions} wurden im YPSDC-Repository vorregistriert. Die YUCT-Koordinationstheorie des Alterns gilt als falsifiziert, wenn nach drei unabhängigen Wiederholungen die experimentellen Ergebnisse für vier der fünf Arten außerhalb der angegebenen Bereiche liegen.
	
	\begin{table}[H]
		\centering
		\caption{Vorregistrierte Vorhersagen (identisch für alle fünf Arten).}
		\label{tab:predictions}
		\begin{tabular}{@{}llll@{}}
			\toprule
			\textbf{Vorhersage} & \textbf{Assay} & \textbf{Erwartet (LL/WT)} & \textbf{Falsifikation wenn} \\
			\midrule
			Mutationsrate       & HPRT           & \(0,86 \pm 0,05\) & LL/WT \(> 0,95\) \\
			ROS                 & MitoSOX        & \(0,80 \pm 0,05\) & LL/WT \(> 0,90\) \\
			Aggregate           & HTT‑Q72-Foci   & \(0,78 \pm 0,05\) & LL/WT \(> 0,90\) \\
			Migrationsgeschwindigkeit & Kratzwundassay & \(1,20 \pm 0,05\) & LL/WT \(< 1,05\) \\
			Replikative Lebensspanne & Kumul. PD     & \(1,5\)–\(2,0\)   & LL/WT \(< 1,3\) \\
			\bottomrule
		\end{tabular}
	\end{table}
	
	\section{Ethische Überlegungen}
	\label{sec:ethics}
	Dieses Protokoll verwendet ausschließlich kommerziell erhältliche Zelllinien. Es werden keine lebenden Tiere benötigt. Die genetischen Konstrukte werden \textit{in vitro} über Standard-Transfektionsreagenzien eingebracht; \textit{in vivo}-Verabreichungsmethoden werden nicht beschrieben. Forscher werden gebeten, Anhang~\(\Sigma\) zu konsultieren, bevor sie diese Arbeit über kultivierte Zellen hinaus ausweiten.
	
	\section{Fazit}
	\label{sec:conclusion}
	Das BCLM‑LL‑MPS-Protokoll liefert einen rigorosen, intern konsistenten Rahmen zur Prüfung der YUCT-Koordinationstheorie des Alterns an fünf evolutionär unterschiedlichen Säugetierarten. Alle Vorhersagen werden aus einem einzigen Satz universeller Konstanten und artspezifischer Codontabellen berechnet; es werden keine freien Parameter \textit{post hoc} angepasst. Ein erfolgreiches Ergebnis würde \(K_{\mathrm{eff}}\) als erste quantitative, artsübergreifende Metrik biologischer Koordination und deren Wiederherstellung als rationale Anti-Aging-Strategie etablieren.
	
	\vspace{0.5cm}
	\noindent\textbf{Daten und Code:} Alle orthologen Sequenzen, Codonnutzungstabellen, berechneten \(\REL\)-Werte und die vorregistrierten Vorhersagen sind verfügbar unter \url{https://github.com/Alexey-Yakushev-YUCT/YPSDC}.
	
	\begin{thebibliography}{99}
		\bibitem{BCLM} A.\,V.\,Yakushev, ``Anhang BCLM: Biologischer Koordinationslagrange,'' in \emph{YUCT}, Zenodo (2026).
		\bibitem{BCLM_LL} A.\,V.\,Yakushev, ``Anhang BCLM‑LL: Experimentelles Protokoll … menschliche Kardiomyozyten,'' in \emph{YUCT}, Zenodo (2026).
		\bibitem{AppendixL} A.\,V.\,Yakushev, ``Anhang L: Fraktale Koordinationsfehlerskalierung,'' in \emph{YUCT}, Zenodo (2026).
		\bibitem{AppendixP} A.\,V.\,Yakushev, ``Anhang P: Temperaturabhängigkeit der Gravitation,'' in \emph{YUCT}, Zenodo (2026).
		\bibitem{AppendixSigma} A.\,V.\,Yakushev, ``Anhang \(\Sigma\): Ethische Imperative und Governance,'' in \emph{YUCT}, Zenodo (2026).
		\bibitem{Bjedov2021} I.~Bjedov \emph{et al.}, ``A single hyper‑accurate mutation in the ribosome extends lifespan in yeast, worms, and flies,'' \emph{Cell Metabolism}, 33, 1054–1070 (2021).
		\bibitem{Kerekes2025} K.~Kerekes \emph{et al.}, ``Codon optimisation and evolutionary pressure in long‑lived mammals,'' \emph{bioRxiv} (2025).
		\bibitem{Morrow2004} G.~Morrow \emph{et al.}, ``Overexpression of … Hsp22 extends \textit{Drosophila} lifespan,'' \emph{FASEB J.}, 18, 598–599 (2004).
		\bibitem{Walker2003} G.\,A.~Walker, G.\,J.~Lithgow, ``Lifespan extension in \textit{C. elegans} by a molecular chaperone,'' \emph{Aging Cell}, 2, 131–139 (2003).
		\bibitem{Takeda2019} M.~Takeda \emph{et al.}, ``Downregulation of \(\alpha\)5 integrin induces cellular senescence,'' \emph{FEBS Lett.}, 593, 1284–1293 (2019).
		\bibitem{codon_usage_db} Y.\,Nakamura \emph{et al.}, Codon usage tabulated … \url{https://www.kazusa.or.jp/codon/}.
		\bibitem{YPSDC_repo} A.\,V.\,Yakushev, YPSDC-Repository, \url{https://github.com/Alexey-Yakushev-YUCT/YPSDC}.
	\end{thebibliography}
	
\end{document}